\chapter{Despre acest suport}
\label{ch:intro}

Fiecare ecran al prezentării calculează pe loc ce arată. Un front-end React
randează cincizeci de ecrane, iar un back-end FastAPI rulează transformata
cerută, la fiecare interacțiune, pe semnalul sau imaginea aleasă de cel care
o urmărește, ori de câte ori pe ecran apare o cifră măsurată: un demo de filtrare aplică un filtru real pe un semnal real, un
demo de compresie rulează o transformată de cosinus discretă și o
transformată wavelet discretă adevărate, pe imagini standard de test.
Rezultatul seamănă mai mult cu un instrument de laborator decât cu un set de
diapozitive: un cursor care schimbă ordinul unui filtru Butterworth și
redesenează răspunsul în frecvență pe loc transmite mai mult decât un
paragraf despre ce înseamnă „ordinul controlează abruptitatea”.

Acest document e gândit pentru cine urmărește prezentarea, live sau din
înregistrare, și vrea ceva scris alături, la care să revină. Aceleași formule,
dar derivate. Aceleași cifre, dar cu metoda din spatele lor. Iar la
electrocardiogramă și la recunoașterea muzicii, mai mult decât arată cele
câteva minute de pe scenă. Capitolele nu urmează neapărat ordinea ecranelor. Sunt grupate pe teme, ca fiecare aplicație practică (compresia imaginilor,
eliminarea zgomotului, ECG, EEG, recunoașterea audio) să se poată citi de
sine stătător, cu teoria strict necesară adusă înainte, nu presărată pe
parcurs.

Prezentarea în sine e un tur ghidat de cincizeci de ecrane, parcurse cu
săgețile tastaturii, grupate în nouă secțiuni vizibile într-o bară laterală
(Figura~\ref{fig:intro-toc}). Fiecare ecran are propriul \emph{deep link} și
poate fi deschis direct, iar secțiunile se pot explora și individual, în
afara turului. Live, la \url{https://ale0204.student-dev.ro/dsp/}.

\begin{figure}[htbp]
\centering
\includegraphics[width=\textwidth]{01-intro-title.png}
\caption{Ecranul de deschidere al turului: cârligul leagă subiectul de o
  cerință de specificație reală. Cinematograful digital impune JPEG2000 în
  locul lui JPEG, prin specificația DCI~\cite{dci2024}.}
\label{fig:intro-title}
\end{figure}

\begin{figure}[htbp]
\centering
\includegraphics[width=\textwidth]{02-intro-toc.png}
\caption{Cuprinsul turului: nouă secțiuni, cu carduri clickabile care sar
  direct la secțiunea aleasă.}
\label{fig:intro-toc}
\end{figure}

\section{Harta: de la ecran la capitol}
\label{sec:harta}

Prezentarea are cincizeci de ecrane, numerotate în bara de jos a aplicației.
Tabelul~\ref{tab:harta} spune, pentru fiecare grup de ecrane, unde se citește
aceeași materie în textul de față. Citit invers, spune ce ecran să deschizi
pentru un capitol anume.

\begin{table}[htbp]
\centering
\caption{Corespondența dintre ecranele turului și capitolele acestui text.}
\label{tab:harta}
\small
\begin{tabular}{@{}llp{5.3cm}@{}}
\toprule
Ecrane & Secțiunea din tur & Unde se citește \\
\midrule
1--2 & Deschidere, cuprins & Capitolul~\ref{ch:intro} \\
3--5 & Transformata Fourier & Capitolul~\ref{ch:fundamente} \\
6--9 & Filtre digitale & Secțiunea~\ref{sec:filtre} \\
10--13 & Convoluția & Capitolul~\ref{ch:fundamente} \\
14--16 & Kernel-uri 2D & Secțiunea~\ref{sec:kernele} \\
17--24 & Transformata wavelet & Capitolul~\ref{ch:wavelets} \\
25--32 & Algoritmul lui Mallat & Secțiunea~\ref{sec:mallat} \\
33--35 & Aplicații, ECG & Capitolul~\ref{ch:ecg} \\
36--37 & EEG & Capitolul~\ref{ch:eeg} \\
38--39 & Amprenta audio & Capitolul~\ref{ch:shazam} \\
40--42 & Denoising & Capitolul~\ref{ch:denoising} \\
43--47 & DCT față de wavelet & Capitolul~\ref{ch:compresie} \\
48--50 & Sinteză, bibliografie & Capitolul~\ref{ch:sinteza} \\
\bottomrule
\end{tabular}
\end{table}

Capitolul~\ref{ch:tehnic} nu are ecrane proprii: descrie cum este construită
și verificată aplicația.
